\documentclass[12pt, a4paper]{article}


\usepackage[utf8]{inputenc}        
\usepackage[english, greek]{babel} 
\usepackage{alphabeta}           
\usepackage[margin=2.5cm]{geometry}
\usepackage{graphicx}              
\usepackage{amsmath, amssymb}      
\usepackage{hyperref}              
\usepackage{setspace}              
\usepackage{fancyhdr}             
\usepackage{listings}
\usepackage[labelfont=bf]{caption}

% \usepackage{libertine}
\usepackage{libertinus}

% Ρυθμίσεις συνδέσμων
\hypersetup{
    colorlinks=false,
    linkcolor=blue,
    urlcolor=cyan,
    pdftitle={Big Data Semester Project Nikolaos Anastasiou},               
    pdfauthor={Nikolaos Anastasiou},             
    pdfsubject={Semester Project - Big Data}, 
    pdfkeywords={big data, spark, hadoop, python}, 
    pdfcreator={XeLaTeX}             
}

\fancyfoot[R]{\thepage}


\onehalfspacing 

\begin{document}


% ==========================================
\begin{titlepage}
    \centering
    \vspace*{1.5cm}
    
    {\huge \textbf{Εθνικό Μετσόβιο Πολυτεχνείο}}\\[0.5cm]
    
    {\Large \textit{Επιστήμη Δεδομένων και Μηχανική Μάθηση}}\\[1.5cm]
    
    {\Large \textbf{Εξαμηνιαία Εργασία}}\\[1cm]

     {\Large \text{Διαχείριση Δεδομένων Μεγάλης Κλίμακας}}\\[1cm]

    \begin{flushleft}
        \textbf{Όνομα / Α.Μ}\\
        Νικόλαος Αναστασίου / 03400283\\[1.5cm]
        

    \end{flushleft}
    
    \vfill
    
    {\large Αθήνα, \today}
    \vspace*{1cm}
\end{titlepage}


\newpage
\begingroup
\hypersetup{hidelinks}
\tableofcontents
\endgroup
\newpage


\phantomsection
\section*{Εισαγωγή}
\addcontentsline{toc}{section}{Εισαγωγή}

Στη παρούσα εργασία υλοποιήθηκαν τα Ζητήματα 1 εως 6 σε περιβάλλον Apache Spark και
Apache Hadoop, ενω ο κώδικας για την απάντηση των Queries γράφτηκε σε Python ακολουθώντας
τα templates που μας δόθηκαν.

Η υλοποίηση έγινε τόσο σε τοπικό επίπεδο για την ανάπτυξη του κώδικα
αλλά και στο απομακρυσμένο cluster που μας δόθηκε πρόσβαση μέσω VPN. Η ανάπτυξη 
έγινε σύμφωνα με τον οδηγό του \href{https://github.com/ikons/bigdata-dsml}{github.com/ikons/bigdata-dsml}.

Αντί να γίνει εγκατάσταση του Spark και του Hadoop στο
WSL των Windows έγινε σε native Ubuntu 24.04.4 LTS καθώς αντί για Windows εχω
Linux. Η μόνη διαφορά είναι η εκδοσή της Java που αντί για Java OpenJDK 11 που ήταν η
προτεινόμενη χρησιμοποιήθηκε η προ εγκατεστημένη δηλαδή η Java OpenJDK 21.

\phantomsection
\section*{Ζητούμενο 1}
\addcontentsline{toc}{section}{Ζητούμενο 1}

Τα δεδομένα που δόθηκαν για την εργασία είναι σε μορφή \texttt{.csv} η οποία αν και είναι
εύκολη στη διαχείριση της και μπορεί να διαβαστεί ανοίγοντας το αρχείο με έναν 
text editor δεν είναι η βέλτιστη για την υποδομή του Spark. Για τον λόγο αυτό
τα δεδομένα μετατράπηκαν στην προτεινόμενη μορφή για Spark που είναι η \texttt{Parquet}.

Υπάρχουν πολλές μορφές δεδομένων οι οποίες είναι καταλληλότερες για Spark έναντι του
απλού \texttt{.csv}, αρχικά είναι η μορφή \texttt{.parquet}, στην οποία θα μετατραπούν
τα δεδομένα της εργασίας, η οποία είναι Columnar δηλαδή αποθηκεύει τα δεδομένα ανα στήλη
και οχι ανα γραμμή. Έτσι αμα ενα query ζητάει μόνο 1 στήλη απο τις 10 τότε απο τον δίσκο
διαβάζεται μονό αυτή η μια στήλη, και με αυτό τον τρόπο πετυχαίνει μεγαλύτερες ταχύτητες 
read, καλύτερη συμπίεση δεδομένων, υποστηρίζει επίσης αλλαγές στη δομή των δεδομένων 
και Predicate Pushdown το οποίο παραλήπτει δεδομένα που δεν χρειάζονται κατά την εκτέλεση 
ενός query αυξάνοντας έτσι ακόμη περισσότερο την ταχύτητα εκτέλεσης. 

Ακόμη ενα format κατάλληλο για Spark είναι το \texttt{ORC (Optimized Row Columnar)},
το οποίο είναι Columnar σαν το \texttt{Parquet}, και έχει πολλές ομοιότητες.
Αναπτύχθηκε για το Apache Hive και λειτουργεί εκεί καλύτερα απο οτι στο Spark. 
Οι διαφορές ανάμεσα στα 2 είναι οτι: Το ORC έχει σαν μικρότερη μονάδα το 
Row Group ενω το Parquet έχει ενα Page 8kb, το ORC έχει Column Pruning το οποίο
προσπερνάει τα columns τα οποία δεν χρειάζονται για κάποιο query ενω το parquet 
έχει column chunks τα οποία είναι ξεχωριστά αποθηκευμένα μέσα στον δίσκο. Επίσης
υπάρχουν διαφορές στο φιλτράρισμα των δεδομένων καθώς στο ORC υπάρχουν φίλτρα για
min/max, bloom filters και offsets ενω στο parquet υπάρχουν αντίστοιχα για 
file, chunk και page.

\newpage

\phantomsection
\section*{Ζητούμενο 2}
\addcontentsline{toc}{section}{Ζητούμενο 2}
Στο ζητούμενο 2 θα δημιουργηθούν 3 υλοποιήσεις για το ίδιο Query όπου στην
πρώτη θα χρησιμοποιεί μονό DataFrame χωρίς User Defined Function (UDF), έπειτα
θα υλοποιηθεί ξανά με UDF και τέλος θα γίνει μια υλοποίηση με Resilient 
Distributed Dataset (RDD). Για ολα τα υπο ερωτήματα τα Crime Data συγχωνευτήκαν σε ενα
ενιαίο αρχείο .csv και .parquet για την ευκολότερη διαχείριση τους.





\end{document}